% !TEX root = ~/Code/TN_LOWE_RECONSTRUCTION/main.tex
\section{Event Simulation}\label{sec:SIMULATION}
In a \ac{LArTPC}, charged particles ionize the argon atoms and an applied electric field causes the liberated electrons to drift toward a set of readout planes. These readout planes are divided into induction and collection planes and, crucially, they are oriented so that the readout direction of the induction plane is rotated with respect to the one of the collection plane; allowing for a 2D reconstruction by matching time--coincident signals. In \ac{DUNE}, the collection plane has its readout wires or strips oriented along constant $z$ values, and there are two induction planes providing useful redundancy in the plane matching to prevent the reconstruction of accidental noise coincidence. For all planes, the arrival times and charges of all signal are recorded, allowing for a 2.5D--reconstruction of the interaction $<t,y,z>$. In addition to ionization signals, \ac{LAr} also produces scintillation light, which can be used for triggering, improved calorimetry, background rejection and timing (reconstructing the drift coordinate $x$). For low--energy events, efficient charge collection and light detection are crucial for a reliable reconstruction. The neutrino event simulation in \ac{DUNE} follows the standard low--energy workflows of the \ac{LArSoft} \ac{MC} framework~\cite{Snider2017wjd}. An example of the energy deposition and charge readout of a low--energy \ac{CC} neutrino interaction in \ac{LAr} can be seen in \cref{fig:EVENT_EDEP}.

\begin{figure}[ht!]
    \centering
    \includegraphics[width=\linewidth]{Images/hd_1x2x6_centralAPA_marley_edep_EDep_event_7501.png}
    \caption[width=1\linewidth]{Example of a simulated \SI{\sim 18}{\mega\electronvolt} neutrino interaction in the \ac{HD} geometry. (left) The event displays the two \ac{TPC} projections of the primary electron energy depositions (orange) and the associated de--excitation gammas (turquoise). The energy deposits from the truth information of the signal propagation are marked with dots and the reconstructed clusters with empty squares. (right) The reconstructed objects are displayed in a 3D representation of the workspace geometry. The represented planes are the central \ac{APA} (green) and the corresponding \ac{CPA}s (orange)\label{fig:EVENT_EDEP}}
\end{figure}

\subsection{Geometries\label{sec:GEOMETRY}}
Due to the high computational load of low--energy productions (including backgrounds), the full scale geometries of the \ac{FD}s are represented in simulation with smaller workspace geometries; designed to provide all possible edge and propagation effects of real neutrino and background interactions but reducing the detector's volume to contain all the signal of a single neutrino event. In this work, we used two geometries: the \ac{HD} workspace geometry (hd\_1x2x6) and the \ac{VD} workspace geometry (vd\_1x8x14\_3view\_30deg). In both cases, the $x$ and $y$ coordinates are set to the same lengths as for the full drift volumes but the $z$ coordinate is shortened roughly $1/4$ and $1/3$ of the full length, respectively. The \ac{HD} workspace features a single central \ac{APA} and two adjacent \ac{CPA}s creating two drift volumes that are designed to emulate the central region of the \ac{HD} detector. On the other hand, the \ac{VD} workspace geometry consists of a single drift volume that represents the upper half of the detector (see \cref{fig:WORKSPACE_GEOMETRY}).

\begin{figure}[ht!]
    \centering
    \includegraphics[width=.49\linewidth]{Images/fd10kt_1x2x6.png}
    \includegraphics[width=.49\linewidth]{Images/fdvd10kt_1x8x14_3view_30deg.png}
    \caption[width=1\linewidth]{(left) The \ac{HD} workspace geometry (hd\_1x2x6\_centralAPA) with a single central \ac{APA} and a drift length of \SI{3.6}{\meter}. (right) The \ac{VD} workspace geometry (vd\_1x8x14\_3view\_30deg).\label{fig:WORKSPACE_GEOMETRY}}
\end{figure}

The \ac{HD} geometry is used in low--energy neutrino studies simulating two background scenarios: one without external backgrounds from the lateral faces of the cavern \ac{APA} (hd\_1x2x6\_centralAPA) and another one with complete external background (hd\_1x2x6\_lateralAPA). Due to the arrangement of the \ac{TPC} planes in the workspace geometry, the latter background scenario is required to implement the lateral cavern flux in the center of the detector (where the \ac{APA} is placed). This creates a scenario where the central region of the workspace actually represents the most external face of the detector's charge readout. In both scenarios, we simulate the full set of radiological backgrounds, among with top, bottom, head, and tail external backgrounds.

For the \ac{VD} geometry, we consider two scenarios: the nominal scenario, which represents the upper half of the \ac{VD} module (vd\_1x8x14\_3view\_30deg), and a shielded scenario that assumes additional reduction of backgrounds in the form of altered fluxes (vd\_1x8x14\_3view\_30deg\_shielded). The latter scenario is likely to represent the lower half of the \ac{VD} module, where shielding is more suitable from and engineering point of view.

\subsection{Neutrino Interaction\label{sec:NEUTRINO_GENERATION}}
The first step in simulating neutrino interactions in DUNE relies on a comprehensive nuclear model of the possible processes. For low--energy events, this is \ac{MARLEY}, an event generator developed to simulate the interactions of low--energy neutrinos in the range of a few to tens of \unit{\mega\electronvolt}. It is particularly suited for simulating interactions relevant to solar and supernova neutrinos in \ac{LAr} detectors. In this work we focus on the \ac{CC} process (although \ac{ES} and \ac{NC} have also been implemented):

\begin{equation}
    \nu_e + {}^{40}\mathrm{Ar} \rightarrow e^- + {}^{40}\mathrm{K}^*
\end{equation}

Upon an electron neutrino capture on an argon nucleus and subsequent electron release, the excitation of the residual \( {}^{40}\mathrm{K}^* \) leads to the emission of de--excitation \(\gamma\)--rays, which can produce observable signals (blips) in the vicinity (\SI{<1}{\meter}) of the main electron track (see \cref{fig:EVENT_EDEP}). \ac{MARLEY} incorporates a gamma--decay schemes to describe these processes and feeds the resulting daughter particles into \ac{LArSoft}, which is used to produce a \ac{MC} input for physics studies. It simulates primary interaction kinematics, the resulting nuclear state and the gamma--cascades from nuclear de--excitation, as well as, other secondary processes (such as neutron or proton emission). In \ac{DUNE}'s physics analyses, the most important signature from these interactions is the already mentioned primary electron track. The resulting energy of these is strongly dependent on the de--excitation scheme of the residual nuclear state (see \cref{fig:NEUTRINO_ELECTRON_CORRELATION}). 

\begin{figure}[ht!]
    \centering
    \includegraphics[width=\linewidth]{Images/hd_1x2x6_centralAPA_marley_True_Electron_Energy.png}
    \caption[width=1\linewidth]{(left) Neutrino--electron energy correlation in a CC interaction as predicted by MARLEY. (right) Gamma energies of the interaction. The solid line splits the energy into two regions with mean values marked with dashed lines.\label{fig:NEUTRINO_ELECTRON_CORRELATION}}
\end{figure}

In \cref{fig:MARLEY_CC_INTERACTION}, the average energy transfer and binding energy of the \ac{CC} interaction are displayed. One can observe how the gamma channel starts dominating the interaction up to \SI{\sim 10}{\mega\electronvolt}, and above that threshold, the main electron takes over the fractional energy. This trend is accompanied by a continuous increase in average kinematic threshold caused by the production of additional secondary particles. The most concerning secondary production is the emission of neutrons, which happens above \SI{\sim15}{\mega\electronvolt}, and leads to a decrease in visible energy that is unlikely to be recovered under realistic reconstruction scenarios. All of these effects combine to a strong energy dependence of the deposited and visible energies, which should be taken into account in our signal calibration. Given that the true neutrino energy reconstruction is our reconstruction goal, the here presented calorimetric algorithms use the number of hits in the primary clusters as a proxy for neutrino energy to encompass the different physics of each energy range.

\begin{figure}[ht!]
    \centering
    \includegraphics[width=.49\linewidth]{Images/hd_1x2x6_TMarleySumK_CC_Mean_Fractional_Channel_Energy.png}
    \includegraphics[width=.49\linewidth]{Images/hd_1x2x6_CC_Kinematic_Threshold.png}
    \caption{\ac{MARLEY}'s computation of \ac{CC} neutrino interaction. (left) Daughter fraction of the original kinetic energy. (right) Average kinematic threshold.\label{fig:MARLEY_CC_INTERACTION}}
\end{figure}

It is important to mention at this point that a typical low--energy \ac{MC} production includes a full propagation of the backgrounds described in~\cref{sec:BACKGROUND_MODEL}. These additional particles are produced in their specific rates and concentrations according to computations out of the scope of this publication. All particle kinematics and detector placement happen in the so--called \emph{gen} stage of the simulation. All produced background particles follow the same workflow chain as the signal particles and contribute to the ultimate smearing and resolution of the reconstructed objects. In order to account for all possible matching issues one full drift--time needs to be simulated before the neutrino makes its interaction (at $t=0$). After this, the simulation last for another full drift--time to allow neutrinos from all regions of the detector to drift and be collected onto the corresponding readout plane. For the purpose of this publication, neutrinos are produced with isotropic positions and directions across the full active volume; and with a flat energy spectrum from \SI{4}{\mega\electronvolt} up to \SI{30}{\mega\electronvolt}.

\subsection{Particle Propagation}
Following the generation of the primary neutrino interaction via the \ac{MARLEY} generator for low--energy solar neutrino events, the resulting final--state particles are propagated through the liquid argon medium using a Geant4 simulation. This so called \emph{g4} stage, models the detailed particle--matter interactions, including energy loss via ionisation, multiple scattering, bremsstrahlung, photon production, and hadronic interactions; according to a physics list tailored for these low--energy processes.

As charged particles traverse the liquid argon, they lose energy primarily through ionization, producing free ionization electrons and excitation of argon atoms, which subsequently decay and emit prompt scintillation photons in the vacuum ultraviolet (VUV) range (\SI{128}{\nano\meter}). These processes are central to the signal generation in a \ac{LArTPC}. For computational efficiency, the simulation does not explicitly track every ionisation electron or VUV photon. Instead, energy loss along particle trajectories is recorded as a set of energy depositions, or \texttt{edeps}, which include position, time, energy, and associated particle identity. Similarly, the number of scintillation photons produced in each energy deposition step is recorded as a secondary output, commonly referred to as \texttt{simphotons}. These \texttt{edeps} and \texttt{simphotons} are then effectively propagated through the \ac{LAr} volume with custom analytical algorithms in a subsequent stage of simulation.

\subsection{Detector Simulation}
In this, often called, \emph{detsim} stage; after the effective propagation of ionization electrons and scintillation photons to the TPC (wires or strips) and PDS (photodetectors) readout elements, the signal formation is modelled. This stage takes as input the energy depositions (\texttt{edeps}) and the associated scintillation photons (\texttt{simphotons}) from the particle propagation step. The detector response to these inputs, produces digitized signals based on the detector's readout electronics and noise characteristics. The two components of the detector response simulation are:

\begin{itemize}
    \item{\textbf{TPC Signal}}.
    Ionisation electrons drift under an applied electric field toward the TPC wire planes. During this drift, the electrons experience diffusion (both longitudinal and transverse), attachment to electronegative impurities (parametrised by the electron lifetime), and charge loss due to recombination. These effects are modelled to modify the total collected charge and the temporal and spatial spread of the signal. \ac{DUNE}s \ac{LArTPC} use three planes (also called ``views'') to detect the drifting charge (1 collection plane and 2 induction planes).

    \begin{enumerate}
        \item\textbf{Collection plane} responds with unipolar signals as the electrons are finally absorbed.
        \item\textbf{Induction planes} respond with bipolar signals as electrons pass by without being collected.
    \end{enumerate}

    The motion of charge induces time--dependent signals (waveforms) on the readsout channels, which are digitized at high frequency. These raw waveforms are then processed using simulated electronic noise based on detector-specific measurements, which is applied to emulate realistic signal-to-noise conditions and electronics shaping. This full set of operations are performed in \ac{LArSoft} using the Wire--Cell toolkit~\cite{Abratenko_2022}. To improve charge extraction, the waveforms undergo a \textbf{deconvolution procedure}, where the electronics response and noise spectrum are inverted in frequency space to recover a cleaner signal representation. The deconvolved waveforms are integrated over time windows to form localized \texttt{hits}, which are stored with time, position, and charge--profile information.

    \item{\textbf{PDS Signal}}.
    Simultaneously, the scintillation photons are detected by the \ac{PDS}, composed of \ac{XA}s. The \texttt{simphotons} from Geant4 are converted into raw signals and integrated in units of \ac{PE} using detector-specific quantum efficiencies and timing models. Similarly to the \ac{TPC} reconstruction, the \ac{PDS} response is also emulated with realistic timing resolution, dark noise, and electronic gain amplification. Simulated waveforms are produced for each optical channel, overlaid with noise, and optionally deconvolved or filtered. The resulting waveforms are processed to extract \emph{optical hits} (integrated photon pulses) with associated time and \ac{PE} information.
\end{itemize}

\subsection{Event Reconstruction}\label{sec:RECONSTRUCTION}
As a final step, after the low--energy simulation steps of generation, propagation and detector simulation of the \ac{MC} workflow. The reconstruction step for these studies includes a basic clustering scheme that applies both to the TPC and PDS hits generating the final TPC clusters and PDS flashes to be used in the analysis. This \emph{reco} stage typically is subdivided into two sub--steps:
\begin{enumerate}
    \item \emph{Reco1}: The TPC \texttt{hits} are disambiguated\footnote{} and grouped according to tight space and time constrains to form \texttt{clusters}, which aim to combine all the charge from localized ionization signals of single particles. The PDS \texttt{ophits} are grouped following the same logic but with a much broader space range to account for the larger spread of scintillation light.
    \item \emph{Reco2}: The \emph{Reco1} \ac{TPC} \texttt{cluster} are further grouped into \texttt{lowe events} containing one \texttt{primary} and any number of found \texttt{adjacent clusters}. For each \texttt{primary cluster}, a flash--matching algorithm is applied to select a single \ac{PDS} \texttt{flash} object to reconstruct the drift--coordinate of the event.
\end{enumerate} 

\begin{figure}[ht!]
    \centering
    \includegraphics[width=\linewidth]{Images/hd_1x2x6_centralAPA_marley_TPC_event_9939.png}
    \caption[width=1\linewidth]{Example of a simulated \SI{\sim 10}{\mega\electronvolt} neutrino interaction in the \ac{HD} geometry. (left) The event displays the two \ac{TPC} projections of the primary electron cluster (orange), the associated de--excitation gammas (turquoise) and uncorrelated background activity (black). The biggest energy deposits from the truth information of the signal propagation are marked with empty squares and the reconstructed clusters with filled squares. (right) The same points are displayed in a 3D representation of the workspace geometry. The represented planes are the central \ac{APA} (green) and the corresponding \ac{CPA}s (orange)\label{fig:EXAMPLE_EVENT}}
\end{figure}

An example of a simulated neutrino interaction in the \ac{HD} geometry is shown in~\cref{fig:EXAMPLE_EVENT}. The event displays the two \ac{TPC} projections of the primary electron cluster (orange), the associated de--excitation gammas (turquoise) and uncorrelated background activity (black). The biggest energy deposits from the truth information of the signal propagation are marked with empty squares and the reconstructed clusters with filled squares. Compared to~\cref{fig:EVENT_EDEP}, the clusters are now reduced to a point--like data structure, which results from the clustering algorithm.

The clustering algorithms developed for this purpose con be found in the following repositories: (\href{https://github.com/DUNE/duneopdet/tree/develop/duneopdet/SolarNuUtils}{duneopdet} \& \href{https://github.com/DUNE/dunereco/tree/develop/dunereco/LowEUtils}{dunereco}); as well as the associated producer modules: (\href{https://github.com/DUNE/duneopdet/blob/develop/duneopdet/OpticalDetector/SolarOpFlash_module.cc}{SolarOpFlash\_module.cc} \& \href{https://github.com/DUNE/dunereco/blob/develop/dunereco/LowEUtils/SolarCluster_module.cc}{SolarCluster\_module.cc}). A custom solar workflow leading to the production of these reconstructed objects has been included in the standard production workflow files: \href{https://github.com/DUNE/dunesw/blob/develop/fcl/dunefd/reco/solar_reco_dune10kt_1x2x6.fcl}{solar\_reco\_dune10kt\_1x2x6.fcl} \& \href{https://github.com/DUNE/dunesw/blob/develop/fcl/dunefdvd/reco/solar_reco_dunevd10kt_1x8x14_3view_30deg.fcl}{solar\_reco\_dunevd10kt\_1x8x14\_3view\_30deg.fcl}.

\subsection{Event Analysis}\label{sec:MC_ANALYSER}
For the specific purpose of the solar neutrino studies an analyzer module (\href{https://github.com/DUNE/duneana/blob/develop/duneana/SolarNuAna/SolarNuAna_module.cc}{SolarNuAna\_module.cc}) has been written to combine these reconstruction objects with the truth information of the event. This module essentially combines all the information needed to perform the calibration and reconstruction of the signal candidates. The biggest caveat from the analysis point of view is that all clusters and flash information is reduced to a point--like data structure. This is, of course, a choice of design in the analyzer and allows for a simpler handling of the data because all event candidates have the same dimension. In future works, this analyzer should be substituted by the more general \href{https://github.com/DUNE/duneana/blob/develop/duneana/LowEAna/LowEAna_module.cc}{LowEAna\_module.cc} which aims to provide the same event reconstruction while keeping the internal structure of the candidate clusters to allow for better calibration and discrimination strategies.